This appendix specifies the SHACL (Shapes Constraint Language) coverage for validating ADL Lite Level~3 (L3) relation blocks against the ADL Core ontology constraints.

\subsection{SHACL Shape Definitions}

ADL Lite L3 relations are validated against a SHACL shape graph with the following constraints:

\begin{enumerate}[label=\textbf{S\arabic*}, leftmargin=3em]
    \item \textbf{Source and target concept existence.} The \texttt{source} and \texttt{target} fields must reference valid \texttt{concept\_id} values that exist in the ADL corpus. Implemented as \texttt{sh:pattern} constraints requiring valid identifiers.
    \item \textbf{Closed predicate set.} The \texttt{relation} field must be one of the registered predicates in \texttt{adl\_core\_ontology.yaml}: \texttt{isomorphic-to}, \texttt{specialisation-of}, \texttt{co-occurs-with}, \texttt{related-to}, \texttt{analogical-to}, \texttt{analogical-transfer}, \texttt{dual-of}, \texttt{fork-of}, \texttt{mitigated-by}, \texttt{indexed-phrase}.
    \item \textbf{Confidence bounds.} The optional \texttt{confidence} field, if present, must satisfy $0.0 \leq \texttt{confidence} \leq 1.0$.
    \item \textbf{Mandatory directionality.} Every relation must have both \texttt{source} and \texttt{target} fields; self-referential relations (\texttt{source} = \texttt{target}) are permitted only for symmetric predicates (\texttt{co-occurs-with}, \texttt{compatible-with}).
    \item \textbf{Mapping type constraint.} The optional \texttt{mapping\_type} field, if present, must be one of \{\texttt{exact}, \texttt{close}, \texttt{broad}, \texttt{narrow}, \texttt{related}\}.
\end{enumerate}

\subsection{SHACL Expressivity Coverage}

Table~\ref{tab:shacl-coverage} maps each L3 relation constraint to the SHACL feature employed and indicates whether it requires SHACL Core (standard) or SHACL-SPARQL (extended) expressivity.

\begin{table}[htbp]
\centering
\caption{SHACL coverage analysis for ADL Lite L3 relation validation. Core = standard SHACL; SPARQL = requires SHACL-SPARQL extension.}
\label{tab:shacl-coverage}
\begin{tabular}{@{}p{3.5cm}p{2.5cm}p{3.5cm}p{3.0cm}@{}}
\toprule
\textbf{Constraint} & \textbf{ADL L3 Field} & \textbf{SHACL Feature} & \textbf{Expressivity} \\
\midrule
Identifier format & \texttt{source}, \texttt{target} & \texttt{sh:pattern} (regex) & Core \\
Predicate membership & \texttt{relation} & \texttt{sh:in} (enumeration) & Core \\
Confidence range & \texttt{confidence} & \texttt{sh:minInclusive}, \texttt{sh:maxInclusive} & Core \\
Directionality & \texttt{source} $\neq$ \texttt{target} & \texttt{sh:sparql} (not-equal check) & SPARQL \\
Mapping type enumeration & \texttt{mapping\_type} & \texttt{sh:in} (enumeration) & Core \\
Concept existence (cross-doc) & \texttt{source}, \texttt{target} & \texttt{sh:sparql} (external lookup) & SPARQL \\
Self-reference symmetry & \texttt{source} = \texttt{target} & \texttt{sh:sparql} (conditional) & SPARQL \\
Timestamp ordering & EventChain linkage & Not expressible in SHACL & N/A \\
\bottomrule
\end{tabular}
\end{table}

\textbf{Coverage summary.} Five of the eight constraints can be expressed in SHACL Core, providing standardisable validation that any SHACL engine can execute without SPARQL support. Three constraints require SHACL-SPARQL: (1) directionality enforcement (source $\neq$ target for asymmetric predicates), (2) cross-document concept existence verification, and (3) conditional self-reference symmetry rules. The EventChain temporal ordering constraint is not expressible in SHACL at all---it requires sequential hash verification, which is a computational rather than a structural constraint. This coverage profile is consistent with SHACL's design intent: structural constraints are declarative (SHACL Core), whereas cross-document and conditional constraints require the expressivity of SPARQL queries.

\subsection{Validation Pipeline}

The SHACL shapes are implemented in \texttt{adl\_lite/shacl\_validation.py} and executed via the \texttt{pyshacl} library against exported relation triples. The shape definitions are generated programmatically from \texttt{adl\_core\_ontology.yaml} rather than being stored as a static \texttt{.ttl} file; they can be serialised on demand using \texttt{adl\_lite/shacl\_validation.py}.
All L3 relation blocks from the experiment corpus (E1--E6) pass SHACL validation with zero constraint violations. The validation pipeline is integrated into the CI workflow via \texttt{scripts/reproduce/reproduce\_shacl.sh}.

\subsection{Implementation}

The SHACL shape graph is generated programmatically by \texttt{adl\_lite/shacl\_validation.py} from the predicate registry in \texttt{adl\_core\_ontology.yaml}. It can be serialised to Turtle format on demand and contains:
\begin{itemize}[leftmargin=2em]
    \item \textbf{ADLCoreRelationShape:} Validates individual L3 relation blocks against constraints S1--S5.
    \item \textbf{ADLChainIntegrityShape:} Validates that exported PROV-O activity sequences preserve the original EventChain temporal ordering (structural check only; cryptographic verification is performed outside SHACL).
\end{itemize}
The shape graph is compatible with standard SHACL engines, including PySHACL, Apache Jena SHACL, and GraphDB SHACL validation.
